% id: disjoint-mutsig
% title: Disjointness and Mutually Singular
% date: 19 Oct. 2024

\documentclass{article}
\usepackage{amsmath, amssymb}

\begin{document}

For an ordered set \((X, \le)\), we say \(x, y \in X\) are disjoint if for any \(z \le x\) and \(z \le y\),
we have that \(z = \perp\). Rémy Degenne on the Lean Zulip observed that, in the case of \(\sigma\)-finite measures,
the notion of mutually singular is equivalent to disjointness. He then asked whether or not this fact generalizes
to arbitrary measures (Rémy's proof used the Radon-Nikodym derivative and consequently relies on
\(\sigma\)-finiteness). In this short note, we provide an elementary proof of this equivalence without
assuming \(\sigma\)-finiteness.

For a measurable space \((X, \Sigma)\) and two measures \(\mu, \nu\) on \((X, \Sigma)\), we define
the measure \(\mu \wedge \nu\) by
\[(\mu \wedge \nu)(S) = \inf \{\mu(S \cap T) + \nu(S \cap T^c) \mid T \in \Sigma\}.\]
It is a standard exercise to check that this is indeed a \textit{bona-fide} measure and that
\(\mu \wedge \nu \le \mu\) and \(\mu \wedge \nu \le \nu\).

Fixing \(\epsilon > 0\) and \(n \in \mathbb{N}\), as \((\mu \wedge \nu)(X) = 0\),
we have by construction that there exists \(S_n^\epsilon \in \Sigma\) such that
\[\mu(S_n^\epsilon) + \nu((S_n^\epsilon)^c) \le \epsilon 2^{-n}.\]
Thus, denoting \(S^\epsilon = \bigcap_{n \in \mathbb{N}} S_n^\epsilon\), we have that
\(\mu(S^\epsilon) = 0\) and
\[\nu((S^\epsilon)^c) = \nu\left(\bigcup_{n \in \mathbb{N}} (S_n^\epsilon)^c\right)
    \le \sum_{n \in \mathbb{N}}\epsilon 2^{-n} = \epsilon.\]
Finally, taking \(S = \bigcup_{n \in \mathbb{N}} S^{1 / n}\),
\(\mu(S) = 0\) as \(S\) is a countable union of \(\mu\)-null-sets while
\[\nu(S^c) = \nu\left(\bigcap_{n \in \mathbb{N}} (S^{1 / n})^c\right) \le \frac{1}{k}\]
for all \(k \in \mathbb{N}\) implying that \(\nu(S^c) = 0\). Thus, \(\mu\) and \(\nu\) are mutually singular
as required.

\end{document}
